\section{Introduction}

Although physics-informed neural networks (\PINN{}s) have shown success on simple partial differential equations (PDEs), they suffer from severe training stagnation and failure when applied to stiff physical regimes \cite{Raissi2019, Cuomo2022review}. Historically, these mesh-free architectures build on early neural PDE solvers \cite{Lagaris1998} and Ritz formulations \cite{E2018deepritz} to optimize governing equations as soft constraints. In this paper, we conduct a systematic empirical investigation mapping the failure landscape of \PINN{}s across diverse architectures, optimization routines, and physical parameters. Our results demonstrate that these training failures are fundamentally structurally determined, contradicting the common assumption that local hyperparameter tuning can resolve stiffness-induced pathologies.

The literature has increasingly documented distinct failure modes that arise when \PINN{}s confront complex physical systems. One prominent limitation is spectral bias, wherein standard multi-layer perceptrons exhibit a strong, inherent tendency to prioritize learning low-frequency components of the target function while struggling disproportionately to resolve high-frequency content \cite{Wang2021gradpathology}. This frequency bias implies that, for PDEs exhibiting multiscale behavior, turbulent cascades, or sharp advective fronts, the network's capacity to uniformly minimize the residual loss across all spatial frequency bands is structurally constrained by the eigenvalue distribution of the associated Neural Tangent Kernel (NTK). Consequently, training may stagnate early in the optimization process, yielding overly smooth solutions that are unable to capture fine-scale dynamics accurately \cite{Wang2021gradpathology}.

A second critical issue identified in recent literature is the phenomenon of gradient pathology, which manifests as severe imbalances in the backpropagated gradients corresponding to different terms of the physics-informed loss function \cite{Wang2021gradpathology}. Analytical and empirical studies indicate that the magnitudes and directions of gradients originating from data fidelity terms, initial or boundary conditions, and PDE residuals often conflict sharply. These conflicting gradient directions suggest that standard gradient descent optimization methods, such as Adam or L-BFGS, struggle to simultaneously satisfy multiple, heterogeneous constraints. This effectively destabilizes the learning process during early training phases and leads to optimization trajectories that stall in sub-optimal local minima characterized by high physical residuals \cite{Wang2021gradpathology, Wang2022causal}.

Additionally, the sampling density and spatial distribution of collocation points have been identified as primary drivers of training failure in both forward and inverse problems \cite{Krishnapriyan2021}. Research indicates that \PINN{}s frequently fail to propagate boundary information correctly into the domain interior when continuous temporal sequences or highly stiff advection regimes are considered. These collocation effects suggest that naive or uniform sampling strategies are fundamentally insufficient to guide the network toward the physically correct solution manifold. In under-sampled regions, this frequently results in non-physical artifacts, artificial numerical dissipation, or catastrophic structural failure, even in scenarios where the global aggregated loss values appear deceptively small \cite{Krishnapriyan2021, Jagtap2020}.

Although these individual pathologies have been meticulously identified and analyzed in isolation, the current body of literature predominantly treats these failure modes as independent phenomena that can be rectified through targeted hyperparameter tuning or isolated architectural interventions. Various strategies—such as adaptive loss weighting schemes, specialized activation functions, domain decomposition methods, or causal training protocols—have been actively proposed to address specific deficiencies \cite{Wang2021gradpathology, Wang2022causal, Jagtap2020}. Even recent high-level surveys and taxonomies of PINN challenges tend to categorize spectral bias, gradient pathologies, and precision errors as fundamentally distinct obstacles, offering ad-hoc mitigation strategies for each in isolation. However, there remains a significant gap in understanding how these mechanisms interact concurrently during the global optimization trajectory. At present, no systematic, joint empirical investigation exists that holistically evaluates the interconnected nature of these failure modes across diverse architectures, optimization strategies, and progressively scaling PDE stiffness regimes. Consequently, it remains unclear whether these documented pathologies are genuinely independent algorithmic artifacts amenable to localized fixes, or if they are instead surface-level symptoms of much deeper, interconnected structural limitations inherent to the continuous physics-informed training paradigm itself.

To address this critical gap in the literature, this paper introduces and empirically evaluates the \textit{Zugzwang} hypothesis—a framing explicitly drawn from the chess concept in which any possible move inevitably worsens a player's position on the board. We hypothesize that \PINN{} failures on stiff PDEs are predominantly structurally determined rather than merely being symptomatic of poor, localized hyperparameter choices. Under this hypothesis, the multiple failure modes that co-occur under high stiffness cannot be resolved by available hyperparameter interventions, not because they form a forced trade-off, but because they arise from independent structural properties of the MLP architecture and composite soft-penalty objective. This perspective suggests that the optimization landscape of \PINN{}s is constrained by fundamental, invariant conditioning bottlenecks. Therefore, attempts to heuristically balance heterogeneous loss terms or dynamically adjust sampling densities may ultimately only trade one manifestation of failure for another. By explicitly framing these intertwined, competing challenges as a structural optimization trap rather than a sequence of isolated algorithmic limitations, we aim to investigate the premise that standard physics-informed continuous architectures inevitably converge toward non-physical local minima when confronted with sufficiently challenging advective or viscous PDE parameters.

To rigorously test the proposed Zugzwang hypothesis, we systematically map the failure landscape of physics-informed neural networks through a series of 24 controlled experiments. The specific contributions of this work are detailed as follows:
\begin{itemize}
    \item We quantitatively characterize the onset of spectral bias across varying advection coefficients and multiple standard activation functions, delineating the strict boundaries of learning failure in advection-dominated regimes (Exps 1--3).
    \item We comprehensively profile the Neural Tangent Kernel (NTK) condition numbers and global Hessian eigenvalue spectra at both initialization and late-stage training, exposing the severe, structural gradient conflicts that persistently stall optimization (Exps 4--6).
    \item We identify and characterize a sharp collocation starvation cliff, directly demonstrating how explicit sampling thresholds dictate the sudden boundary between successful convergence and abrupt, catastrophic structural failure (Exp 8).
    \item We empirically demonstrate the presence of severe temporal extrapolation failure, revealing how these networks structurally fail to generalize beyond the immediate training domain despite exhibiting deceptively negligible interior loss characteristics (Exp 17).
    \item Two-dimensional failure phase space mapping across $(\beta, \Ncol{})$, $(\nu, \Ncol{})$, and $(\beta, \text{width})$ parameter planes, revealing a discontinuous phase wall where PINNs transition instantaneously from trivially solvable to structurally impossible, with no intermediate scaling regime amenable to parametric fitting (Exp.~25).
    \item We perform rigorous variance decomposition analysis on the Burgers' equation to systematically isolate the relative impact of underlying structural factors versus tuning hyperparameter choices on overall model stability (Exp.~27).
    \item Conservation law audit revealing that gradient pathology (FM3) and collocation starvation (FM5) achieve nearly identical $\Ltwo$ errors ($0.418$ vs $0.416$) yet differ by $57\times$ in mass conservation violation, despite near-identical energy conservation violation (Experiment~26, Section~\ref{sec:exp26}).
\end{itemize}

The remainder of this paper is organized as follows. Section \ref{sec:methods} formally outlines our foundational experimental setup, PDE problem definitions, and core evaluation metrics. Subsequent sections are dedicated to detailing the specific empirical findings associated with distinct, yet interacting, failure mechanisms: the impact of spectral barriers (Section \ref{sec:part1}), the geometry of optimization traps (Section \ref{sec:part2}), and the manifestation of data illusions (Section \ref{sec:part3}). In Section \ref{sec:part4}, we explicitly evaluate existing literature mitigation strategies within the contextual bounds of the Zugzwang hypothesis. Further rigorous analyses of network phase space dynamics, inherent conservation property violations, and global variance decomposition are presented in Sections \ref{sec:part5}, \ref{sec:part6}, and \ref{sec:part7}, respectively. Finally, Section \ref{sec:conclusion} synthesizes our findings, concludes the study, and highlights critical implications for future directions in scientific machine learning architectures.
